\section{线性逆问题的解析解：伪逆引导}

当前向映射 $\mathcal{A}$ 是线性函数时（这涵盖了去模糊、超分辨率、修复、CT 重建等大量实际应用），我们可以得到似然分数的解析表达式。

\subsection{线性前向模型的设定}

假设 $\mathcal{A}$ 是线性的，即 $y = A x_0 + n$，其中 $A$ 是矩阵，$n \sim \mathcal{N}(0, \sigma_y^2 I)$ 是观测噪声。此时：

$$
p(y | x_0) = \mathcal{N}(y; A x_0, \sigma_y^2 I)
$$

\begin{knowledgebox}{常见线性逆问题中的 $A$ 矩阵}
\begin{itemize}
    \item \textbf{高斯去模糊}：$A$ 是高斯卷积核对应的矩阵。
    \item \textbf{超分辨率}：$A$ 是下采样矩阵。
    \item \textbf{图像修复}：$A$ 是掩码对角矩阵（已知区域为 1，未知区域为 0）。
    \item \textbf{CT/MRI 重建}：$A$ 是投影/傅里叶采样矩阵。
\end{itemize}
\end{knowledgebox}

\subsection{高斯共轭先验的推导}

利用高斯分布的共轭性质，当先验 $p(x_0 | x_t)$ 和似然 $p(y|x_0)$ 都是高斯分布时，边际分布 $p(y | x_t)$ 也是高斯分布。具体地：

\textbf{先验}（高斯近似后）：
$$
p(x_0 | x_t) \approx \mathcal{N}(\hat{x}_0, r_t^2 I)
$$

\textbf{似然}：
$$
p(y | x_0) = \mathcal{N}(A x_0, \sigma_y^2 I)
$$

\textbf{边际}（对 $x_0$ 积分后）：
$$
p(y | x_t) = \int p(y | x_0) p(x_0 | x_t) \, dx_0 = \mathcal{N}(A \hat{x}_0, \; r_t^2 A A^\top + \sigma_y^2 I)
$$

\begin{importantbox}{伪逆引导公式（Pseudo-Inverse Guidance）}
基于上述高斯边际分布，似然分数可以解析计算：
$$
\nabla_{x_t} \log p(y | x_t) \approx -\frac{1}{2} \nabla_{x_t} \left\| y - A \hat{x}_0 \right\|^2_{(r_t^2 A A^\top + \sigma_y^2 I)^{-1}}
$$
其中 $\hat{x}_0$ 是通过 Tweedie 公式从 $x_t$ 计算得到的去噪估计。

当观测无噪声（$\sigma_y = 0$）时，公式简化为需要计算 $A$ 的\textbf{伪逆}（pseudo-inverse），因此被称为伪逆引导。
\end{importantbox}

\begin{warningbox}{伪逆引导仅适用于线性 $\mathcal{A}$}
此解析公式的推导依赖于 $\mathcal{A}$ 的线性性。对于非线性前向映射（如涉及神经网络的渲染、复杂物理模拟等），需要使用更一般（但近似更粗糙）的方法。
\end{warningbox}

\subsection{本章小结}

对于线性逆问题，通过高斯共轭先验的性质，我们可以得到似然分数的解析表达式。该方法需要计算前向映射矩阵的伪逆，因此称为伪逆引导。它为去模糊、超分辨率、修复等经典图像处理任务提供了理论上较为严谨的解法。


